\documentclass[11pt]{article}
\usepackage[a4paper,margin=24mm]{geometry}
\usepackage[T1]{fontenc}
\usepackage{amsmath,amssymb,amsthm,mathtools}
\usepackage[hidelinks,hypertexnames=false]{hyperref}
\setlength{\emergencystretch}{3em}
\newcommand{\R}{\mathbb R}\newcommand{\C}{\mathbb C}
\title{Independent verification of the literal theta--prolate map\newline and its complete divided-quotient jets}
\author{}\date{20 September 2026}
\begin{document}\maketitle

\section{Scope, source and corrections}
This is an independent calculation of PT1--16 in the companion
\emph{The literal prolate--theta source map: constants, truncation tails and original divisor jets}.
The actual author LaTeX of Connes--Consani--Moscovici \cite{CCM},
from the \emph{Outlook} section through \emph{The missing steps},
was read for this check: equations \texttt{ku}, \texttt{emap},
\texttt{xih}, \texttt{ktoh}, and the complete proofs of Lemmas
\texttt{hermfact}, \texttt{meixnerlem}, \texttt{hermfact1}.
The source is the original file \texttt{mc2arXiv.tex} of
arXiv:2511.22755v1. No PDF, numerical figure, or original
Meixner--Sch\"afke or Fuchs paper was read for this verification.
Their estimates remain attributed through the inspected CCM argument.

The scalar seed, its Mellin transform, all three truncation terms,
and the inverse spectral coordinate in PT1--16 are correct.
Three repairs were reported to the author of the companion:
the input function in PT8 must be measurable as well as bounded;
the compact convergence proved by PT10 is on the open critical strip,
not on arbitrary compact subsets of $\C$;
and the draft PT5 contained a control character in place of
the command beginning its second fraction. No original author
source was edited. The quotient functions below are entire individually;
that statement does not enlarge the region in which convergence has
been proved.

\section{The literal seed and every transform factor}
Use exactly
\[
 \begin{gathered}
 H(x)=(\pi^2x^4-\tfrac32\pi x^2)e^{-\pi x^2},\quad
 \phi_*=4H,\quad \Theta\phi(x)=2\sum_{n\ge1}\phi(nx),\\
 f_0=8\sum_{n\ge1}H(nx),\quad K(u)=u^{1/2}f_0(u)/8.
 \end{gathered}
\]
For $\Re s>0$, the substitution $t=\pi x^2$ gives
\[
 \int_0^\infty x^{s+r-1}e^{-\pi x^2}\,dx
 =\frac{\Gamma((s+r)/2)}{2\pi^{(s+r)/2}}.
\]
Thus, retaining both terms in $H$,
\begin{equation}\tag{IC1}
 \mathcal MH(s)=\pi^{-s/2}\Gamma(s/2)
 \left[\tfrac12\tfrac s2(\tfrac s2+1)-\tfrac34\tfrac s2\right]
 =\frac{s(s-1)}8\pi^{-s/2}\Gamma(s/2).
\end{equation}
For $\Re s>1$ the integral of the absolute sum is bounded by
$\zeta(\Re s)\int_0^\infty|H(x)|x^{\Re s-1}\,dx$.
It follows that
\[
 \mathcal Mf_0(s)=s(s-1)\pi^{-s/2}\Gamma(s/2)\zeta(s)=g(s)=2\xi(s).
\]

The two Hermite functions displayed in CCM give directly
\begin{equation}\tag{IC2}
 H=\frac{\sqrt3}{2^{11/4}}h_4-\frac3{2^{17/4}}h_0,
 \qquad \|H\|_{L^2(\R)}^2=\frac3{2^{11/2}}+\frac9{2^{17/2}}
 =\frac{33}{2^{17/2}}.
\end{equation}
The norm can also be computed by expanding $H^2$ and using
$\int_{\R}x^{2j}e^{-2\pi x^2}\,dx
=\Gamma(j+1/2)/(2\pi)^{j+1/2}$.
The Gaussian transform with kernel $e^{2\pi ixy}$ is
$e^{-\pi y^2}$. Multiplication by $x^r$ transforms to
$(2\pi i)^{-r}\partial_y^r$. Applying this rule for $r=2,4$
to the displayed polynomial proves $\mathbb FH=H$.
In particular $H(0)=\int_\R H=0$. Poisson summation now gives
\begin{equation}\tag{IC3}
 \sum_{n\ge1}H(nu)=u^{-1}\sum_{n\ge1}H(n/u),
 \quad f_0(u)=u^{-1}f_0(1/u),\quad K(u)=K(1/u).
\end{equation}
Gaussian decay and IC3 prove rapid decay at both endpoints,
including every scaling derivative. Hence the Mellin integrals
of $f_0$ and the multiplicative Fourier integrals of $K$ are entire.
The half-plane calculation extends to every $s$.
With the precise transform convention
$\widehat K(z)=\int_0^\infty K(u)u^{-iz}\,du/u$, one obtains
\begin{equation}\tag{IC4}
 \widehat K(z)=\tfrac18g(1/2-iz)
 =\tfrac14\xi(1/2+iz)=\tfrac14\Xi(z).
\end{equation}
The sign change uses IC3, equivalently $\xi(1-s)=\xi(s)$.
Thus the literal seed in CCM's equation \texttt{ku} has transform
$\Xi/4$ under its stated conventions. The seed $4H$ has transform
$\Xi$; this scalar change cannot be dropped in masses or Gram matrices.

The map $UF=u^{1/2}F$ from $L^2(du)$ to $L^2(du/u)$ is unitary,
because both squared norms equal $\int|F|^2du$. Its inverse is
the multiplication by $u^{-1/2}$. Direct differentiation yields
\begin{equation}\tag{IC5}
 K=Uf_0/8,\quad \|K\|^2=\|f_0\|^2/64,\quad
 U(-u\partial_u)U^{-1}=\tfrac12-u\partial_u,
 \quad -iu\partial_u=i(U D U^{-1}-\tfrac12).
\end{equation}
These identities retain the original centre, phase and mass.
Also $\mathcal E(\phi_*)=4K=U\Theta\phi_*/2$ by the two displayed
sum definitions, without any change of the factor-two trace.

\section{The complete closed-substrip estimate}
Let $\lambda\ge1$ and let $A_\lambda$ be bounded and measurable
on $[-\lambda,\lambda]$, extended by zero outside that interval.
Set
\[
 d_\lambda=\sup_{|x|\le\lambda}|A_\lambda(x)-H(x)|,\qquad
 K_\lambda(u)=\mathbf1_{[\lambda^{-1},\lambda]}(u)\sqrt u
       \sum_{1\le n\le\lambda/u}A_\lambda(nu).
\]
This transform is entire because its defining measurable integrand
has compact support away from $0$.
For $u$ in that support the exact error is
\begin{equation}\tag{IC6}
 K_\lambda(u)-K(u)
 =\sqrt u\sum_{n\le\lambda/u}(A_\lambda(nu)-H(nu))
   -\sqrt u\sum_{n>\lambda/u}H(nu).
\end{equation}
The second sum does not vanish: for example $H(n)>0$ for each
integer $n\ge1$. Consequently taking $A_\lambda=H$ on its interval
does not justify dropping this term.

For $v\ge1$ put
\[
 B(v)=(\pi^2v^4+\tfrac32\pi v^2)e^{-\pi v^2},\quad
 T(v)=\int_v^\infty B(x)\,dx.
\]
Here $|H(v)|\le B(v)$ and
\[
 B'(v)=\pi v e^{-\pi v^2}(-2\pi^2v^4+\pi v^2+3)<0.
\]
Indeed $-2y^2+y+3<0$ for $y\ge\pi>3/2$.
Separating the first omitted lattice point and bounding subsequent
points by the integral of this decreasing function gives
\begin{equation}\tag{IC7}
 \sum_{n>\lambda/u}|H(nu)|\le B(\lambda)+T(\lambda)/u.
\end{equation}
Integration by parts, or differentiation followed by the value at
infinity, gives exactly
\begin{equation}\tag{IC8}
 T(v)=e^{-\pi v^2}(\tfrac\pi2v^3+\tfrac32v)
       +\tfrac34\operatorname{erfc}(\sqrt\pi v),\qquad
 J_\nu(v)=\int_v^\infty x^\nu e^{-\pi x^2}\,dx
 =\frac{\Gamma((\nu+1)/2,\pi v^2)}{2\pi^{(\nu+1)/2}}.
\end{equation}

Let $0\le a<1/2$ and $\alpha=\Im z$ with $|\alpha|\le a$.
The two integrals required by IC6--7 are
\[
 \int_{1/\lambda}^{\lambda}u^{\alpha-3/2}\,du
 =\frac{\lambda^{1/2-\alpha}-\lambda^{\alpha-1/2}}{1/2-\alpha},
 \quad
 \int_{1/\lambda}^{\lambda}u^{\alpha-1/2}\,du
 =\frac{\lambda^{1/2+\alpha}-\lambda^{-1/2-\alpha}}{1/2+\alpha}.
\]
Each is at most
$I_a=(\lambda^{1/2+a}-\lambda^{-1/2-a})/(1/2+a)$.
To see this, use the increasing function
$b\mapsto2\sinh(b\log\lambda)/b$ for $b>0$; its derivative
has the sign of $t\cosh t-\sinh t$, nonnegative by differentiation
from $t=0$.

For $u\ge\lambda$, $|K(u)|\le\sqrt u B(u)+u^{-1/2}T(u)$.
Reflection IC3 turns the missing lower endpoint into an upper
endpoint with exponent $-\alpha$. Both endpoint errors are thus
at most
\[
 2\int_\lambda^\infty
     \big[u^{a-1/2}B(u)+u^{a-3/2}T(u)\big]\,du.
\]
Positive Fubini evaluates the latter $T$ integral as
\[
 \int_\lambda^\infty B(v)
       \frac{\lambda^{a-1/2}-v^{a-1/2}}{1/2-a}\,dv
 \le\frac{\lambda^{a-1/2}T(\lambda)}{1/2-a}.
\]
Combining these equalities proves the complete bound
\begin{equation}\tag{IC9}
 \begin{split}
 \sup_{|\Im z|\le a}|\widehat K_\lambda(z)-\Xi(z)/4|
 \le{}&\lambda d_\lambda I_a+[B(\lambda)+T(\lambda)]I_a\\
 &+2\left[\pi^2J_{a+7/2}(\lambda)
 +\tfrac32\pi J_{a+3/2}(\lambda)
 +\frac{\lambda^{a-1/2}T(\lambda)}{1/2-a}\right].
 \end{split}
\end{equation}
All constants and exponents match PT9--12. There is no dependence
on $\Re z$. Substitution of the CCM-attributed prolate estimate
$d_\lambda=O(\lambda^{-2})$ makes its first term
$O(\lambda^{-1/2+a})$. The two Gaussian terms also tend to zero.
The actual discrepancy is retained; no approximation to a minimum
Weil eigenvector has been assumed.

\section{The exact divided-quotient jets}
Let $h(s)=\prod_\rho(s-\rho)^{m_\rho}$ for the actual finite
packet of distinct roots and complete positive multiplicities,
and $d=\sum m_\rho$. Put $b_\rho=h/(s-\rho)^{m_\rho}$ and define
the divided derivatives and local reciprocal coefficients by
\[
 f_{\rho,n}=F^{(n)}(\rho)/n!,\quad
 \beta_{\rho,a}=[t^a]b_\rho(\rho+t)^{-1},\quad
 c_{\rho,j}=\sum_{a=0}^j\beta_{\rho,a}f_{\rho,j-a}.
\]
Taylor multiplication gives
\begin{equation}\tag{IC10}
 P=\mathcal I_hF=\sum_\rho b_\rho(s)
                   \sum_{j=0}^{m_\rho-1}c_{\rho,j}(s-\rho)^j.
\end{equation}
At $\rho$ the other summands vanish to order $m_\rho$, while
the local summand agrees with $F$ to that order. Every summand
has degree less than $d$. A difference of two such polynomials
would be divisible by $h$ and have degree less than $d$, hence
would vanish. This proves interpolation and uniqueness.
Therefore $\mathcal D_hF=(F-P)/h$ is entire for entire $F$.

Write $p_{\rho,n}=P^{(n)}(\rho)/n!$, zero for $n\ge d$.
Taylor division now gives the complete quotient-jet formula
\begin{equation}\tag{IC11}
 \frac{(\mathcal D_hF)^{(\ell)}(\rho)}{\ell!}
 =\sum_{a=0}^{\ell}\beta_{\rho,a}
        (f_{\rho,m_\rho+\ell-a}-p_{\rho,m_\rho+\ell-a}).
\end{equation}
Indeed the numerator has Taylor coefficients zero below
$m_\rho$, and division by $(s-\rho)^{m_\rho}$ shifts those
indices by precisely $m_\rho$. Multiplication by $1/b_\rho$
then yields IC11.

For a formula containing only input divided jets, the identity
$P/h=\sum_\sigma\sum_{j<m_\sigma}c_{\sigma,j}
(s-\sigma)^{j-m_\sigma}$ yields
\begin{equation}\tag{IC12}
 \begin{split}
 \frac{(\mathcal D_hF)^{(\ell)}(\rho)}{\ell!}
 ={}&\sum_{a=0}^{m_\rho+\ell}\beta_{\rho,a}f_{\rho,m_\rho+\ell-a}\\
 &-(-1)^\ell\sum_{\sigma\ne\rho}\sum_{j=0}^{m_\sigma-1}
 c_{\sigma,j}\binom{m_\sigma-j+\ell-1}{\ell}
 (\rho-\sigma)^{-(m_\sigma-j+\ell)}.
 \end{split}
\end{equation}
The coefficient follows by expanding each negative integral
power at $\rho$. The reciprocal coefficients themselves are
the following finite expressions, retaining every root difference:
\begin{equation}\tag{IC13}
 \beta_{\rho,a}=(-1)^a\prod_{\sigma\ne\rho}(\rho-\sigma)^{-m_\sigma}
 \sum_{\substack{k_\sigma\ge0\\\sum k_\sigma=a}}
 \prod_{\sigma\ne\rho}\binom{m_\sigma+k_\sigma-1}{k_\sigma}
                    (\rho-\sigma)^{-k_\sigma}.
\end{equation}
This is the coefficient of $t^a$ in the product of the binomial
series $(1+t/(\rho-\sigma))^{-m_\sigma}$, each multiplied by
its unchanged constant $(\rho-\sigma)^{-m_\sigma}$.
For a single root the empty product is $1$ at $a=0$ and the
empty constrained sum is zero at $a>0$.

IC11--12 show that the quotient jet of order $\ell$ requires
input derivatives through $m_\rho+\ell$ at $\rho$ and only
the interpolation jets at other roots. The entire original quotient
jet at $\rho$, of orders $\ell<m_\rho$, therefore uses derivatives
through $2m_\rho-1$. This order is necessary in the unrestricted
holomorphic class: for $F=h(s)(s-\rho)^\ell$, every interpolation
jet is zero, yet $\mathcal D_hF=(s-\rho)^\ell$ has divided jet
$1$ at order $\ell$. Every derivative of $F$ at $\rho$ below
$m_\rho+\ell$ is zero, and the next is
$(m_\rho+\ell)!b_\rho(\rho)\ne0$.

For a finite quantitative bound choose a circle of radius $r_\sigma$
at every root and put
\[
 S_\sigma=\sup_{|s-\sigma|=r_\sigma}|F(s)|,
 \quad A_{\sigma,n}=\sum_{a=0}^n|\beta_{\sigma,a}|r_\sigma^{-(n-a)}.
\]
Cauchy's integral formula gives
$|f_{\sigma,j}|\le S_\sigma r_\sigma^{-j}$ and
$|c_{\sigma,n}|\le S_\sigma A_{\sigma,n}$. Consequently IC12
proves the explicit finite estimate
\begin{equation}\tag{IC14}
 \begin{split}
 \left|\frac{(\mathcal D_hF)^{(\ell)}(\rho)}{\ell!}\right|
 \le{}&S_\rho A_{\rho,m_\rho+\ell}\\
 &+\sum_{\sigma\ne\rho}S_\sigma\sum_{j<m_\sigma}
 A_{\sigma,j}\binom{m_\sigma-j+\ell-1}{\ell}
 |\rho-\sigma|^{-(m_\sigma-j+\ell)}.
 \end{split}
\end{equation}
These circles need only lie in the domain of $F$; the reciprocal
coefficients are exact local algebraic constants, so another root
inside one of the circles does not invalidate this estimate.

Finally set $G_\lambda(s)=8\widehat K_\lambda(i(s-1/2))$.
Since $-i\,i(s-1/2)=s-1/2$, IC4 gives the exact target $g(s)$.
Also $\Im(i(s-1/2))=\Re s-1/2$. Thus IC9 bounds
$G_\lambda-g$ by eight times its right side on every closed
substrip $|\Re s-1/2|\le a<1/2$. Apply IC14 to this error.
The circles for a fixed packet can all be chosen inside such
a larger closed substrip, proving convergence of every original
quotient jet whenever the displayed transform error tends to zero.
For a compact subset of $0<\Re s<1$, take a bounded contour
inside that strip enclosing both the compact set and packet,
and avoiding the roots. On the contour $|h|$ has a positive
minimum, IC10 controls the interpolant of the error, and division
by $h$ controls its divided error. The maximum principle extends
the bound to the enclosed compact set.

The exact physical unit is still $j_h(g/h)$; its approximant is
$j_h(\mathcal D_hG_\lambda)$ with the higher derivatives specified
above. The sharp interval cutoff may create endpoint jumps.
None of these entire-function, strip, or fixed-packet statements
places that inverse Mellin source in the original strong smooth
space, proves an increasing-degree moment estimate, or proves
that the approximants have only real spectral zeros.

\section{Exact finite verification}
The companion script \texttt{verify\_theta\_prolate.py} checks
the literal Hermite coefficients and Gaussian squared norm;
the Mellin factor, Fourier invariance, generator and inverse-coordinate
signs; the full tail derivatives and incomplete-gamma constants;
the two interior power integrals; and IC10--13 on four explicitly
artificial repeated-root packets, including nonreal roots.
Those packets are illustrative polynomial data, not zeta zeros.
There are 266 exact identities and six deliberate-failure controls.
The latter reject the missing factor four, the missing factor eight,
the missing squared-mass factor 64, a reversed inverse-coordinate
phase before imposing symmetry, an omitted erfc term, and inference
of a quotient jet from only vanishing lower numerator jets.
These finite checks accompany the analytic proofs above; they do
not replace their convergence or domain arguments.

\begin{thebibliography}{9}
\bibitem{CCM} A. Connes, C. Consani and H. Moscovici,
\emph{Zeta Spectral Triples}, arXiv:2511.22755v1
(27 November 2025), \url{https://arxiv.org/abs/2511.22755v1}.
Original author file \texttt{mc2arXiv.tex}; actual inspected scope
is the \emph{Outlook} through \emph{The missing steps}, with the
equations and proofs specified in Section 1. Earlier original
sources cited there are not claimed independently read here.
\end{thebibliography}
\end{document}
